\bookchapter{บทที่ 44: ลงมือจากศูนย์สู่หนึ่ง—โปรเจกต์เต็มรูปแบบครบทุกขั้นตอน}

\chapterintro{ทฤษฎีจบแล้ว ถึงเวลาลงมือทำจริง บทนี้ไม่ใช่บทเรียนแบบ “พิมพ์ตาม” ธรรมดา แต่เป็นวิธีคิดลงมือทำแบบ “คิดยังไง ถามยังไง ตรวจยังไง” เราจะใช้สามโปรเจกต์ที่ยากต่างกัน เพื่อโชว์กระบวนการคิดเต็มรูปแบบจากความต้องการถึงผลงาน ไม่ว่าคุณใช้เครื่องมือไหน ขั้นตอนชุดนี้ใช้ได้ทั้งหมด}

\booksection{โปรเจกต์ที่หนึ่ง: เว็บเครื่องมือส่วนตัวสุดมินิมอล (ระดับเริ่มต้น)}

\booksubsection{คิดเรื่องความต้องการ}

เป้าหมาย: สร้างเว็บเครื่องมือส่วนตัวที่แสดงทางเข้าเครื่องมือเล็ก ๆ ใช้งานได้จริง 3-4 ตัว เป็นโปรเจกต์ฝั่งฟรอนต์ล้วน ไม่ต้องมีแบ็กเอนด์และฐานข้อมูล จุดสำคัญไม่ใช่เครื่องมือซับซ้อนแค่ไหน แต่คือการลองกระบวนการเต็มรูปแบบ “จากความต้องการถึงผลงาน”

ตามองค์ประกอบความต้องการห้าข้อ: Who—ใช้ส่วนตัว What—เว็บหน้าเดียวแสดงทางเข้าเครื่องมือ Why—รวมเครื่องมือที่ใช้บ่อยไว้หน้าเดียวให้เข้าถึงง่าย How—ไฟล์เดียว HTML+CSS+JS ไม่ใช้เฟรมเวิร์ก Constraint—เสร็จใน 30 นาที ไม่ต้องดีพลอย

\booksubsection{สั่งงาน}

ตามหลัก Think Before Coding คำสั่งของคุณควรมีข้อมูลพอให้ AI เข้าใจเป้าหมาย แต่ไม่ต้องกำหนดรายละเอียดการทำงานทุกจุด:

\begin{lstlisting}
ช่วยสร้างหน้าเว็บเครื่องมือส่วนตัว โดยมีเงื่อนไข:
1. HTML ไฟล์เดียว (ฝัง CSS และ JS) ไม่ใช้เฟรมเวิร์ก
2. ธีมมืด สไตล์มินิมอลสมัยใหม่
3. แสดงการ์ดเครื่องมือ 4 ตัว:
   - เครื่องมือสรุปข้อความ
   - เครื่องมือนับคำ
   - เครื่องมือเลือกสี
   - เครื่องมือสร้าง QR Code
4. แต่ละการ์ดมีไอคอน ชื่อ และคำอธิบายสั้น ๆ
5. การ์ดมีแอนิเมชันตอนชี้เมาส์
6. เรสปอนซีฟ ใช้บนมือถือได้
7. ด้านบนมีแถบชื่อเรื่อง
ขอไฟล์ HTML เต็มมาให้เลย
\end{lstlisting}

\booksubsection{ตรวจรับและวนซ้ำ}

เมื่อ AI ส่งโค้ดมา ตรวจรับด้วยการตรวจสามชั้นของการวนซ้ำ ชั้นแรกตรวจฟังก์ชัน: เปิดไฟล์ในเบราว์เซอร์ได้ไหม? การ์ดครบ 4 ใบแสดงครบไหม? อนิเมชันตอนชี้เมาส์ปกติไหม? ชั้นสองตรวจขอบเขต: เลย์เอาต์ปกติในขนาดจอมือถือไหม? หน้าต่างแคบมาก การ์ดขึ้นบรรทัดใหม่ไหม? ชั้นสามตรวจสถาปัตยกรรม: โค้ดอ่านได้ไหม? CSS รวมอยู่ในแท็ก style เดียวไหม? มีส่วนซ้ำซ้อนชัด ๆ ไหม?

ถ้าอยากทำฟังก์ชันของการ์ดเครื่องมือตัวใดด้วย (ไม่ใช่แค่ทางเข้า) ค่อย ๆ เพิ่มทีละอย่าง—เพิ่มฟังก์ชันเครื่องมือหนึ่งตัว ตรวจแล้วค่อยเพิ่มตัวถัดไป นี่คือการใช้กลยุทธ์ “ก้าวเล็กวิ่งไว” อย่างเป็นรูปธรรม

\booksection{โปรเจกต์ที่สอง: ระบบบันทึกรายรับรายจ่ายแบบฟูลสแตก (ระดับก้าวหน้า)}

\booksubsection{คิดความต้องการและเขียน PRD}

เป้าหมาย: แอปบันทึกรายรับรายจ่ายแบบฟูลสแตกที่มีฟรอนต์ แบ็กเอนด์ และฐานข้อมูล โปรเจกต์นี้ซับซ้อนกว่าอันแรกมาก ต้องใช้วิธีวิศวกรรมความต้องการจากบทที่ 40

เริ่มจากเขียน PRD ตามเทมเพลตในหัวข้อ 3 ของบทที่ 40 เขียนภาพรวมโปรเจกต์ รายการฟีเจอร์ แผนเทคนิค โมเดลข้อมูล และแผนพัฒนา การเลือกแนวทางเทคนิคต้องยึดหลัก Simplicity First—ฟรอนต์ใช้ React แต่ไม่เอา Redux หรือ state management ซับซ้อน แบ็กเอนด์ใช้ Node.js แต่ไม่เอาไมโครเซอร์วิส ฐานข้อมูลใช้ SQLite แทน PostgreSQL (ช่วงพัฒนา)

แผนพัฒนาแบ่งเป็น 5 ขั้น: Step 1 โครง API แบ็กเอนด์ (เซิร์ฟเวอร์ Express + SQLite + อินเทอร์เฟซ CRUD); Step 2 หน้าพื้นฐานฝั่งฟรอนต์ (รายการบันทึก + ฟอร์มเพิ่มรายการ); Step 3 เชื่อมฟรอนต์-แบ็ก (เรียก API + แสดงข้อมูล); Step 4 ฟีเจอร์หมวดหมู่และตัวกรอง; Step 5 สรุปรายเดือนและกราฟง่าย ๆ ทุกขั้นจบต้องเป็น “ของที่ใช้ได้”

\booksubsection{ตั้งไฟล์ความจำโปรเจกต์}

ก่อนเขียนโค้ด สร้างไฟล์ความจำโปรเจกต์ก่อน นี่คือการปฏิบัติหัวใจของ context engineering—ทำให้ AI รู้พื้นหลังโปรเจกต์ตั้งแต่ต้นบทสนทนาทุกครั้ง

\begin{lstlisting}
# โปรเจกต์: แอปบันทึกรายรับรายจ่ายแบบเต็มสแตก
## ภาพรวม
แอปบันทึกรายรับรายจ่าย รองรับการบันทึก หมวดหมู่ และสรุปรายเดือน
## เทคโนโลยี
- Frontend: React + TypeScript + Vite
- Backend: Node.js + Express + TypeScript
- ฐานข้อมูล: SQLite (ผ่าน better-sqlite3)
- ไทป์ใช้ร่วมกันทั้งหน้าและหลัง
## ข้อตกลง API
- เส้นทางพื้นฐาน: /api
- รูปแบบคำตอบ: { success: boolean, data: any, error: string }
- สไตล์ RESTful
## โมเดลข้อมูล
Record { id, amount, type, category, note, date }
## ข้อตกลงการเขียนโค้ด
- TypeScript โหมดเข้มงวด
- ย่อหน้า 2 ช่องว่าง
- ฟังก์ชันมีคอมเมนต์ JSDoc
- คอมเมนต์และข้อความคอมมิตเป็นภาษาจีน
\end{lstlisting}

\booksubsection{ทำทีละขั้น}

ตามแผนพัฒนา แต่ละขั้นใช้บทสนทนาแยกกัน นี่คือตัวอย่างคำสั่งของขั้นแรก:

\begin{lstlisting}
ต้องการโครง API ฝั่งแบ็กเอนด์
ข้อกำหนด:
1. เซิร์ฟเวอร์ Express พอร์ต 3001
2. ฐานข้อมูล SQLite ชื่อตาราง records
3. ฟิลด์ตาราง records:
   id (TEXT PRIMARY KEY)
   amount (REAL)
   type (TEXT: 'income'|'expense')
   category (TEXT)
   note (TEXT)
   date (TEXT: YYYY-MM-DD)
4. API CRUD:
   GET /api/records - ดึงรายการทั้งหมด
   POST /api/records - เพิ่มรายการ
   PUT /api/records/:id - แก้ไขรายการ
   DELETE /api/records/:id - ลบรายการ
5. Middleware จัดการข้อผิดพลาด
ขอโครงโค้ดแบ็กเอนด์แบบเต็มมาให้หน่อย
\end{lstlisting}

สังเกตจุดเด่นของคำสั่งนี้: อธิบายแค่ขั้นปัจจุบัน (โครงแบ็กเอนด์) ไม่ยัดเรื่องฟรอนต์; ระบุฟิลด์และอินเทอร์เฟซชัด (ให้ AI มีข้อจำกัดพอ) แต่วิธีทำจริงปล่อยให้ AI ตัดสิน; และยึดหลัก Goal-Driven Execution—บอก “ทำอะไร” ไม่ใช่ “ทำยังไง”

เมื่อ AI ส่งโค้ดมา ตรวจทันที: รันเซิร์ฟเวอร์ ทดสอบทุกอินเทอร์เฟซด้วย curl หรือ Postman ว่าทำงานปกติไหม ยืนยันแล้วเปิดบทสนทนาใหม่ไปทำ Step 2 (หน้าพื้นฐานฝั่งฟรอนต์) ทุกขั้นยึดวงจร “สั่งงาน-ตรวจ-เปิดบทสนทนาใหม่”

\booksection{โปรเจกต์ที่สาม: แอปเสริมพลังด้วย AI (ระดับสูง)}

\booksubsection{ระดมไอเดียโปรเจกต์}

เป้าหมาย: แอปที่เชื่อม API โมเดลใหญ่—เช่น “ผู้ช่วยเรียน AI” ผู้ใช้ป้อนเนื้อหาที่เรียน AI สรุปใจความ เชื่อมโยงความรู้ที่เกี่ยวข้อง และแนะนำแหล่งเรียนรู้ จุดพิเศษคือ: แอปของคุณเองต้องเรียกใช้ AI คุณต้องจัดการเรื่องใหม่ ๆ อย่างการเรียก API, prompt engineering และการจัดการข้อผิดพลาด

บทเรียนสำคัญที่สุดของโปรเจกต์นี้คือ: เมื่อคุณ “ใช้ AI เขียนโค้ดที่ไปเรียกใช้ AI” คุณต้องใส่ใจ prompt engineering เป็นพิเศษ prompt ที่แอปของคุณใช้เรียก API โมเดลใหญ่ก็ต้องยึดหลักที่คู่มือนี้พูดถึง—ชัดเจน กระชับ ขับเคลื่อนด้วยเป้าหมาย

\booksubsection{จุดเทคนิคสำคัญ}

โปรเจกต์นี้มีจุดเทคนิคที่สองโปรเจกต์ก่อนหน้าไม่มี:

\begin{itemize}
  \item การจัดการคีย์ API: ห้ามเขียนคีย์แปะในโค้ดเด็ดขาด จัดการด้วยตัวแปรสภาพแวดล้อม;
  \item การจัดการการตอบแบบสตรีม: คำตอบของโมเดลใหญ่มักมาแบบสตรีม ต้องแสดงทีละตัวอักษร;
  \item การควบคุมการใช้ Token: โมเดลใหญ่คิดเงินตาม token ต้องคุมปริมาณการใช้ในระดับแอป;
  \item การจัดการเทมเพลต Prompt: ฟีเจอร์ต่างกันใช้เทมเพลต prompt ต่างกัน ต้องจัดการอย่างเป็นระบบ
\end{itemize}

จุดเทคนิคพวกนี้ให้ AI ช่วยทำได้หมด แต่คุณต้องเข้าใจ “ทำไม” ของแต่ละจุด—ทำไมคีย์เขียนแปะในโค้ดไม่ได้? ทำไมต้องมีการตอบแบบสตรีม? ทำไมการใช้ token ต้องควบคุม? พอเข้าใจ “ทำไม” คุณถึงตัดสินได้ว่าแนวทางที่ AI เสนอสมเหตุสมผลไหม

\booksection{เช็กลิสต์ขั้นตอนปฏิบัติแบบใช้ร่วมได้}

ไม่ว่าโปรเจกต์เล็กใหญ่ นี่คือเช็กลิสต์ขั้นตอนเขียนโปรแกรมด้วย AI แบบใช้ร่วมได้ทั่วไป ทุกครั้งเริ่มโปรเจกต์ใหม่ ไล่เช็กลิสต์นี้รอบหนึ่ง ให้มั่นใจว่าไม่พลาดขั้นตอนสำคัญ

\begin{center}
\begin{tabularx}{\linewidth}{|X|X|X|}
\hline
\textbf{ระยะ} & \textbf{ทำอะไร} & \textbf{หลักที่เกี่ยวข้อง} \tabularnewline
\hline
1. คิดเรื่องความต้องการ & ไล่องค์ประกอบความต้องการห้าข้อ คิดให้ชัดว่าจะทำอะไร & Think Before Coding \tabularnewline
2. แยกความต้องการ & แยกความต้องการใหญ่เป็นขั้นเล็กที่ตรวจได้อิสระ & Simplicity First \tabularnewline
3. เขียน PRD & จัดเป็นเอกสารมีโครงสร้างไว้อ้างอิง & Goal-Driven \tabularnewline
4. ตั้งไฟล์ความจำ & สร้าง CLAUDE.md หรือไฟล์ลักษณะเดียวกัน & Harness Engineering \tabularnewline
5. ทำทีละขั้น & หนึ่งขั้นหนึ่งบทสนทนา สั่งงาน-ตรวจ-ขั้นถัดไป & Surgical Changes \tabularnewline
6. ทดสอบตรวจสอบ & ทดสอบความถูกต้องของฟังก์ชันทันทีที่แต่ละขั้นเสร็จ & Iterate \& Verify \tabularnewline
7. รีวิวโค้ด & ตรวจความสมเหตุสมผลของสถาปัตยกรรม ความปลอดภัย และการดูแลรักษา & Iterate \& Verify \tabularnewline
8. วนซ้ำปรับปรุง & ปรับปรุงตามผลทดสอบและการรีวิว & Iterate \& Verify \tabularnewline
\hline
\end{tabularx}
\end{center}

\begin{center}\small\color{gray}เช็กลิสต์ขั้นตอนเขียนโปรแกรมด้วย AI แบบใช้ร่วมได้\end{center}

\begin{tcbpractice}{ลงมือทำ: ทำ "โปรเจกต์ที่หนึ่ง" ให้เสร็จในสัปดาห์นี้}
ใช้ “เว็บเครื่องมือส่วนตัว” ในหัวข้อแรกของบทที่ 44 เป็นการบ้านเขียนโปรแกรม AI ชิ้นแรกของคุณ: ใช้ 5 นาทีไล่ห้าองค์ประกอบ ส่งคำสั่งให้ AI (เครื่องมือไหนก็ได้) แล้วตรวจรับด้วยการตรวจสามชั้น เสร็จใน 30 นาที—ผลงานเขียนโปรแกรมด้วย AI ชิ้นแรกของคุณก็ถือกำเนิดขึ้น ทำเสร็จเก็บไว้ดี ๆ นี่คือของที่ระลึกของการเปลี่ยนจาก “ผู้อ่าน” เป็น “ผู้ลงมือทำ”

\end{tcbpractice}

\begin{tcbsummary}{สรุปบทนี้}
\begin{itemize}
  \item โปรเจกต์ที่หนึ่ง (เริ่มต้น): เว็บเครื่องมือฝั่งฟรอนต์ล้วน—ลองกระบวนการครบวงจร เสร็จได้ใน 30 นาที
  \item โปรเจกต์ที่สอง (ก้าวหน้า): ระบบบันทึกรายรับรายจ่ายแบบฟูลสแตก—PRD + ไฟล์ความจำ + ทำทีละขั้น 5 ขั้น ทุกขั้น "ใช้ได้"
  \item โปรเจกต์ที่สาม (สูง): ผู้ช่วยเรียน AI—ใช้ AI เขียนโค้ดที่เรียกใช้ AI ระวังเรื่อง prompt engineering
  \item ขั้นตอนรวม 8 ขั้น: คิดความต้องการ แยกงาน PRD ไฟล์ความจำ ทำทีละขั้น ทดสอบ รีวิว วนซ้ำ
  \item ทุกขั้นยึดวงจร "สั่งงาน-ตรวจ-เปิดบทสนทนาใหม่"
\end{itemize}

\end{tcbsummary}
